% ===============================================================
\chapter{Kripke, Naming and Necessity}
% ===============================================================

\textbf{Saul Kripke}: Snyder's Noether. The God in his eyes. Resurrector of metaphysics.

\section{Deconstructing descriptivism}

Types of descriptivism that Kripke identified:

\begin{definition}\leavevmode
\begin{enumerate}[label=\roman*.]
  \item \emph{Fregean/Russell descriptivism}: names are synonymous with definite descriptions.
  \item \emph{Cluster theory}: names are synonymous with clusters of definite descriptions.
\end{enumerate}
\end{definition}

\begin{argument}[Kripke against descriptivism]
``Desc.'' refers to both Fregean descriptions and cluster descriptions.
\begin{premises}
  \item[P1)] If descriptivism were correct, then for any name $a$ there would be a $d \in$ Desc.\ such that the statement `$a$ is $d$' is (i) analytic, (ii) knowable a priori, (iii) necessarily true.
  \item[P2)] There exists an $a$ such that there does not exist a $d$ for which the above three properties are fulfilled.
  \item[C)] Descriptivism is incorrect and does not hold.
\end{premises}
\end{argument}

\begin{definition}[The three properties]\leavevmode
\begin{enumerate}[label=\roman*.]
  \item \emph{Analyticity}: true in virtue of meaning alone, for example definitions.
  \begin{itemize}
    \item Contrast, \emph{syntheticity}: true in virtue of how the world is.
  \end{itemize}
  \item \emph{Knowable a priori}: knowable independently of experience, for example mathematical statements.
  \begin{itemize}
    \item Contrast, \emph{a posteriori}: knowledge dependent on experience.
  \end{itemize}
  \item \emph{Necessarily true}: true in all possible worlds, for example the law of self-identity.
  \begin{itemize}
    \item Contrast, \emph{contingent}: may or may not be true.
  \end{itemize}
\end{enumerate}
\end{definition}

\begin{note}
Analyticity entails the other two properties.
\end{note}

\section{The metaphysical argument against descriptivism}

\subsection{Modal logic}

To be very honest with you, the slides do an excellent job here, and I am going to leave them to do their job, except for one note. You might notice that the relation of accessibility should be a little more sharply defined, considering that the two modal operators rely on accessing different worlds. The best answer I found was from this guy on StackExchange, which works for the purpose of these notes:
\begin{enumerate}
  \item Assume we are using a variable domain approach, as David Lewis does to preserve his counterparts. If you have defined a proper model, you know that $R(w_1, w_2)$ if in your model the pair $(w_1, w_2)$ is in the extension of $R$. If you do not have a model, you may be able to infer this from the presumed properties of $R$; for example, from $R(w_0, w_1)$ and $R(w_0, w_2)$ you can infer $R(w_1, w_2)$ if $R$ is Euclidean. You define this relation in the way that characterizes the modality in question. You need to have an idea what properties the respective modality has in the first place.
  \item The last two sentences are a little beyond the scope of this course and I am going to leave it alone now. For now, assume that within a certain possibility, if it is not possible to contain a relation within those laws or a certain modality, it is simply inaccessible.
\end{enumerate}

\begin{slide}
The modern notion of necessity comes from \textbf{modal logic}: the logic of possibility and necessity. Its semantics is due to Kripke, who introduces two key notions:
\begin{enumerate}
  \item \textbf{possible worlds}: There are innumerable ways the world could have been: dinosaurs may have survived, Miley Cyrus could have been President, etc. For each way the world could be, there is a possible world in which things are that way. Formally, we can think of these as just points.
  \item \textbf{accessibility}: Possible worlds are related to one another through the relation of accessibility. Intuitively, a possible world $w_2$ is accessible from world $w_1$ if $w_2$ is possible given how things are in $w_1$. There are various formal properties we can impose on the accessibility relation (reflexivity, symmetry, transitivity, etc), resulting in different \emph{systems} of modal logic.
\end{enumerate}
We can think of different modal systems as modeling different kinds of possibility:
\begin{itemize}
  \item \textbf{epistemic possibility}: what's possible given what we know
  \item \textbf{deontic possibility}: what's possible given the laws of morality
  \item \textbf{physical possibility}: what's possible given the laws of physics
  \item \textbf{logical possibility}: what's possible given the laws of logic
  \item \textbf{metaphysical possibility}: what's possible given the ``laws of metaphysics''
\end{itemize}
\end{slide}

\begin{slide}
\textbf{Necessity and Possibility.} Modal logic introduces two operators: `$\Box$' (necessity) and `$\Diamond$' (possibility). Their semantics is given as follows, for any proposition $p$:
\begin{enumerate}
  \item `$\Box p$' is true at $w$ iff $p$ is true at all worlds accessible from $w$
  \item `$\Diamond p$' is true at $w$ iff $p$ is true at least one world accessible from $w$
\end{enumerate}
Given these definitions, we can distinguish between:
\begin{enumerate}
  \item \textbf{necessary truths}: statements true in all accessible worlds
  \item \textbf{necessary falsehoods}: statements false in all accessible worlds
  \item \textbf{contingencies}: statements true in some but not all accessible worlds
\end{enumerate}
\end{slide}

\subsection{The argument}

\begin{argument}[Metaphysical argument]\leavevmode
\begin{premises}
  \item[P0)] Assume `Aristotle' and `the best student of Plato' are synonymous.
  \item[P1)] If descriptivism were correct, then `Aristotle is the best student of Plato' would be analytic, and thus necessarily true.
  \item[P2)] The above sentence is contingent.
  \item[C)] The descriptive theory is wrong.
\end{premises}
\end{argument}

\section{The semantic argument against descriptivism}

\subsection{Rigid designation}

\begin{definition}\leavevmode
\begin{enumerate}[label=\roman*.]
  \item \emph{Designator}: a coverall term for singular terms: expressions whose semantic function is to refer to particular objects, such as names and definite descriptions.
  \item \emph{Rigid designator}: a designator that refers to the same thing in every possible world in which that thing exists.
\end{enumerate}
\end{definition}

\paragraph{Rigidity implies necessity of identity.}
Suppose `$\alpha$' and `$\beta$' are rigid designators, and suppose `$\alpha = \beta$' is true. Then in the actual world, `$\alpha$' and `$\beta$' refer to the same thing, say $y$. But since `$\alpha$' and `$\beta$' are rigid, they refer to $y$ in all possible worlds (in which $y$ exists). But then in every possible world, `$\alpha = \beta$' is true. In other words, `$\alpha = \beta$' is necessarily true.

\subsection{The argument}

\begin{argument}[Semantic argument]\leavevmode
\begin{premises}
  \item[P1)] If descriptivism were correct, then `Aristotle' and `the best student of Plato' would both be rigid designators, referring to Aristotle in all possible worlds.
  \item[P2)] `Aristotle', but not `the best student of Plato', is a rigid designator.
  \item[C)] Thus, descriptivism cannot be correct.
\end{premises}
\end{argument}

\section{The epistemic argument against descriptivism}

(This guy fucking hated descriptivism.)

\begin{argument}[Epistemic argument]\leavevmode
\begin{premises}
  \item[P1)] If descriptivism were correct, then identity statements of the form `$\alpha$ is the $\delta$' would be analytic, and thus knowable a priori. Hence:
  \begin{enumerate}[label=\roman*.]
    \item Speakers knowing the meaning of `$\alpha$' should also know a specific description $\delta$ which distinguishes the referent of `$\alpha$' from everything else.
    \item Speakers should not be able to use `$\alpha$' to refer to $\alpha$ if something other than $\alpha$ is the actual referent of `the $\delta$'.
  \end{enumerate}
\end{premises}
\end{argument}

\begin{example}[Counterexample]
Many people know that Richard Feynman and Murray Gell-Mann are famous physicists. But this is probably all they know. Thus, they don't know a specific $\delta$ which would distinguish Feynman from Gell-Mann. Yet they can refer to both using their names.
\end{example}

\section{Necessity, a prioricity and analyticity}

All the above arguments are different types of arguments based on analyticity (semantic), necessity (metaphysical) and a prioricity (epistemic). There are examples where one property can be absent (analyticity still entails the other two, though):
\begin{enumerate}[label=\roman*.]
  \item \emph{Necessary a posteriori}: knowable only empirically, while still being necessarily true, which implies a lack of a prioricity. For example, ``water is H$_2$O'', which, while true, is only knowable through experience.
  \item \emph{Contingent a priori}: true in some but not all worlds, but knowable independently of experience. For example, ``The length of the standard meter stick $S$ is one meter'', where the length is a non-rigid designator.
\end{enumerate}

\section[Direct reference and causal reference theory]{Direct reference and causal reference theory\\ (post-Kripkeans)}

\begin{definition}[Direct reference theory]
The meaning of a name is its referent. That is, at least for names, there is no distinction between meaning and reference.
\end{definition}

\paragraph{Semantics vs metasemantics.}
\begin{enumerate}
  \item \emph{Semantic theories}: a theory telling us what any name means (descriptivism, direct reference theory).
  \item \emph{Metasemantic theories}: a theory telling us why the reference-relation holds between a name and its referent, if it holds at all. The causal theory is hinted at for this one.
\end{enumerate}

\begin{definition}[Causal theory of reference]\leavevmode
\begin{enumerate}[label=\roman*.]
  \item For a name `$\alpha$', an \emph{initial baptism} happens, in one of two ways:
  \begin{enumerate}
    \item via \emph{ostension}: by pointing directly at an object and dubbing it `$\alpha$';
    \item via \emph{description}: by giving a description `$\delta$' and declaring that the referent of `$\alpha$' is to be whatever uniquely satisfies `$\delta$'.
  \end{enumerate}
  \item The initial baptism fixes the referent of `$\alpha$'.
  \item The name is then passed from the initial dubber to other speakers, who use the name with the intention of referring to the same object.
  \item This establishes a causal link between the referent and subsequent users of the name.
  \item The name continues being passed in this manner, from speaker to speaker, establishing a causal-historical link between speakers and the referent.
  \item It is in virtue of there being a causal chain linking any user of the name and the referent that speakers manage to successfully refer to that thing using `$\alpha$'. Without that link, no reference would be achieved.
\end{enumerate}
\end{definition}

Direct reference theory is just the Referential Theory of Meaning applied to names. Descriptivism has been torn apart and we are back to a very interesting square one.
